\documentclass[../main.tex]{subfiles}
\begin{document}
\clearpage
\appendix
\section{高速検出器チャンネルマップライブラリ}
\label{sec:appendix}
マップアルゴリズムの比較検討

データ収集や解析において、生データに記録されているヒットデータと検出器の情報を対応づける、つまり検出器チャンネルマッピングは必要不可欠である。この章では主にハイレベルフィルター内での使用に採用した検出器チャンネルマッピングについて述べる。

C++ STLのコンテナには二種類のアルゴリズムのマップが提供されている。それは\lstinline|std::unordered_map|と\lstinline|std::map|である。前者はマップハッシュテーブルを用いたマップであり、後者は二分探索木を用いたマップである。前者は計算量$\mathcal{O}(1)$であり、後者は計算量$\mathcal{O}(\log n)$である。しかしながらそのどちらも採用せず、独自の検出器チャンネルマップを実装し、ライブラリとして公開した。
\begin{itemize}
  \item フロントエンドの読み出し基板のIPアドレス, 読み出しチャンネル番号のセットが、どの検出器のチャンネルに対応するか分かること
  \item 高速に動作すること
\end{itemize}
\subsection{二分探索木を用いたマップアルゴリズムの検討}
\subsection{計算量$\mathcal{O}(1)$のマップアルゴリズムの比較}
\subsubsection{}
\subsection{CUDAでも使用可能な同様のアルゴリズムのマップライブラリ}
\end{document}
